\documentclass[10pt,a4paper]{amsart}
\usepackage{amsmath,amssymb,amsthm,graphicx}
% Retained preamble for the Stacks Project chapter {CH} (shared project preamble).
% The chapter file pinned-original-source.tex only carries \input{{preamble}} on its first line,
% so no preamble text lives inside the chapter itself.  The Stacks Project uses one single shared
% preamble for every chapter: the snapshot stages it three times (sources/stacks-chow/preamble.tex,
% sources/stacks-etale-cohomology/preamble.tex and focused-seeds/AGP/sources/stacks-snapshot-20261011/
% preamble.tex) and those copies are identical in size.  The block below is copied mechanically,
% line for line, from sources/stacks-chow/preamble.tex lines 14-18 and 154-217: the Xy-pic setup and
% the native theorem declarations together with the mathematical macros.  Nothing was added, removed
% or redefined; the wrapper class and the amsmath/amssymb/amsthm/graphicx base packages are supplied
% by the assembler.
% For commutative diagrams we use Xy-pic
% Chapter: stacks-cohomology.tex -- Cohomology of Sheaves.
\usepackage[all]{xy}

% We use 2cell for 2-commutative diagrams.
\xyoption{2cell}

% Theorem environments.
%
\theoremstyle{plain}
\newtheorem{theorem}[subsection]{Theorem}
\newtheorem{proposition}[subsection]{Proposition}
\newtheorem{lemma}[subsection]{Lemma}

\theoremstyle{definition}
\newtheorem{definition}[subsection]{Definition}
\newtheorem{example}[subsection]{Example}
\newtheorem{exercise}[subsection]{Exercise}
\newtheorem{situation}[subsection]{Situation}

\theoremstyle{remark}
\newtheorem{remark}[subsection]{Remark}
\newtheorem{remarks}[subsection]{Remarks}

\numberwithin{equation}{subsection}

% Macros
%
\def\lim{\mathop{\mathrm{lim}}\nolimits}
\def\colim{\mathop{\mathrm{colim}}\nolimits}
\def\Spec{\mathop{\mathrm{Spec}}}
\def\Hom{\mathop{\mathrm{Hom}}\nolimits}
\def\Ext{\mathop{\mathrm{Ext}}\nolimits}
\def\SheafHom{\mathop{\mathcal{H}\!\mathit{om}}\nolimits}
\def\SheafExt{\mathop{\mathcal{E}\!\mathit{xt}}\nolimits}
\def\Sch{\mathit{Sch}}
\def\Mor{\mathop{\mathrm{Mor}}\nolimits}
\def\Ob{\mathop{\mathrm{Ob}}\nolimits}
\def\Sh{\mathop{\mathit{Sh}}\nolimits}
\def\NL{\mathop{N\!L}\nolimits}
\def\CH{\mathop{\mathrm{CH}}\nolimits}
\def\proetale{{pro\text{-}\acute{e}tale}}
\def\etale{{\acute{e}tale}}
\def\QCoh{\mathit{QCoh}}
\def\Ker{\mathop{\mathrm{Ker}}}
\def\Im{\mathop{\mathrm{Im}}}
\def\Coker{\mathop{\mathrm{Coker}}}
\def\Coim{\mathop{\mathrm{Coim}}}

% Boxtimes
%
\DeclareMathSymbol{\boxtimes}{\mathbin}{AMSa}{"02}

%
% Macros for moduli stacks/spaces
%
\def\QCohstack{\mathcal{QC}\!\mathit{oh}}
\def\Cohstack{\mathcal{C}\!\mathit{oh}}
\def\Spacesstack{\mathcal{S}\!\mathit{paces}}
\def\Quotfunctor{\mathrm{Quot}}
\def\Hilbfunctor{\mathrm{Hilb}}
\def\Curvesstack{\mathcal{C}\!\mathit{urves}}
\def\Polarizedstack{\mathcal{P}\!\mathit{olarized}}
\def\Complexesstack{\mathcal{C}\!\mathit{omplexes}}
% \Pic is the operator that assigns to X its picard group, usage \Pic(X)
% \Picardstack_{X/B} denotes the Picard stack of X over B
% \Picardfunctor_{X/B} denotes the Picard functor of X over B
\def\Pic{\mathop{\mathrm{Pic}}\nolimits}
\def\Picardstack{\mathcal{P}\!\mathit{ic}}
\def\Picardfunctor{\mathrm{Pic}}
\def\Deformationcategory{\mathcal{D}\!\mathit{ef}}

% --- appended 2026-10-12 (Chief): the project own definitions of the reference/slogan/history
% display environments, copied mechanically from the shared Stacks preamble lines 8-12
% (\usepackage{verbatim} plus the three \newenvironment lines, which the chapter preamble keeps
% verbatim from the project preamble). Nothing added, removed or redefined: these environments
% carry no mathematical content and the project itself typesets them as comments. Without these
% three lines a counted statement whose own text uses them fails with Environment slogan undefined.
\usepackage{verbatim}
\newenvironment{reference}{\comment}{\endcomment}
\newenvironment{slogan}{\comment}{\endcomment}
\newenvironment{history}{\comment}{\endcomment}
\usepackage{cleveref}
\crefname{originallemma}{original lemma}{original lemmas}
\Crefname{originallemma}{Original lemma}{Original lemmas}
\makeatletter
\newlabel{derived-lemma-higher-derived-functors}{{015B}{source}{}{sourcepointer-1}{}}
\newlabel{derived-lemma-higher-derived-functors@cref}{{[originallemma][0][]015B}{[1][0][]source}{}{}{}}
\makeatother
\crefname{originallemma}{original lemma}{original lemmas}
\Crefname{originallemma}{Original lemma}{Original lemmas}
\makeatletter
\newlabel{lemma-cech-cohomology-delta-functor-presheaves}{{01EK}{source}{}{sourcepointer-2}{}}
\newlabel{lemma-cech-cohomology-delta-functor-presheaves@cref}{{[originallemma][0][]01EK}{[1][0][]source}{}{}{}}
\makeatother
\crefname{originallemma}{original lemma}{original lemmas}
\Crefname{originallemma}{Original lemma}{Original lemmas}
\makeatletter
\newlabel{lemma-injective-trivial-cech}{{01EP}{source}{}{sourcepointer-3}{}}
\newlabel{lemma-injective-trivial-cech@cref}{{[originallemma][0][]01EP}{[1][0][]source}{}{}{}}
\makeatother
\crefname{originallemma}{original lemma}{original lemmas}
\Crefname{originallemma}{Original lemma}{Original lemmas}
\makeatletter
\newlabel{lemma-ses-cech-h1}{{01EU}{source}{}{sourcepointer-4}{}}
\newlabel{lemma-ses-cech-h1@cref}{{[originallemma][0][]01EU}{[1][0][]source}{}{}{}}
\makeatother
\title{Vanishing of higher Cech cohomology for all open coverings implies vanishing of cohomology}
\author{The Stacks Project Authors}
\date{stacks-project:cohomology.tex:local-snapshot-20261011; GNU FDL 1.2 or later}
\begin{document}
\maketitle
\noindent\textit{Editorial note.} All delivered mathematical text is the author's own, copied from the pinned Stacks Project chapter file cohomology.tex; no statement or proof step was rewritten, reordered or removed.\par\medskip
\noindent\textit{Editorial note.} This card was re-admitted after a gate audit: it had been excluded only because the chapter's retained preamble lacks the slogan environment. The author's slogan and the Xy-pic diagram of the proof are both part of the counted spans; the card carries a per-card preamble copy formed from the chapter retained preamble plus the project's own three display environments copied verbatim from sources/stacks-chow/preamble.tex lines 7-12.\par\medskip
\noindent\textit{Editorial note.} The counted spans carry 4 navigation stubs and no bibliography citation.\par\medskip
\noindent\textit{Editorial note.} The label lemma-cech-vanish is defined only in cohomology.tex, so no cross-chapter restatement had to be weighed for this card. The variant of this lemma stated on a basis of opens is a different statement counted by its own card.\par\medskip
\noindent\textit{Allowed standard prerequisites (editorial).}\begin{itemize}
\item Injective modules have trivial higher Cech cohomology. Same chapter, Lemma lemma-injective-trivial-cech, invoked at proof lines 1652-1653. Supplies the vanishing for the injective used in the induction.
\item The hypothesis upgrades a short exact sequence of sheaves to a short exact sequence of presheaves. Same chapter, Lemma lemma-ses-cech-h1, invoked at proof lines 1659-1660. Supplies the exactness of the sequence of presheaves and hence the Cech vanishing of Q.
\item Cech cohomology is a cohomological functor on presheaves of modules. Same chapter, Lemma lemma-cech-cohomology-delta-functor-presheaves, invoked at proof lines 1661-1663. Supplies the long exact sequence used in the induction.
\item Higher cohomology of an injective module vanishes. Derived Categories, Lemma derived-lemma-higher-derived-functors, invoked at proof lines 1683-1685. Supplies H\textasciicircum{}n(U, I) = 0 for n > 0.
\end{itemize}\medskip
\ifdefined\hypertarget\hypertarget{sourcepointer-1}{}\fi
\noindent\textit{Source reference.} Original lemma 015B: sources/stacks-chow/tags.txt, source lines 1502--1502. Pointer outside the counted span, invoked once for context; the load-bearing argument of the counted statement is the author's proof reproduced in full in the proof span.\par\medskip
\ifdefined\hypertarget\hypertarget{sourcepointer-2}{}\fi
\noindent\textit{Source reference.} Original lemma 01EK: sources/stacks-cohomology/pinned-original-source.tex, source lines 1066--1066. Pointer outside the counted span, invoked once for context; the load-bearing argument of the counted statement is the author's proof reproduced in full in the proof span.\par\medskip
\ifdefined\hypertarget\hypertarget{sourcepointer-3}{}\fi
\noindent\textit{Source reference.} Original lemma 01EP: sources/stacks-cohomology/pinned-original-source.tex, source lines 1408--1408. Pointer outside the counted span, invoked once for context; the load-bearing argument of the counted statement is the author's proof reproduced in full in the proof span.\par\medskip
\ifdefined\hypertarget\hypertarget{sourcepointer-4}{}\fi
\noindent\textit{Source reference.} Original lemma 01EU: sources/stacks-cohomology/pinned-original-source.tex, source lines 1593--1593. Pointer outside the counted span, invoked once for context; the load-bearing argument of the counted statement is the author's proof reproduced in full in the proof span.\par\medskip
\begin{lemma}
\label{lemma-cech-vanish}
\begin{slogan}
If higher {\v C}ech cohomology of an abelian sheaf vanishes for all open covers,
then higher cohomology vanishes.
\end{slogan}
Let $X$ be a ringed space.
Let $\mathcal{F}$ be an $\mathcal{O}_X$-module such that
$$
\check{H}^p(\mathcal{U}, \mathcal{F}) = 0
$$
for all $p > 0$ and any open covering $\mathcal{U} : U = \bigcup_{i \in I} U_i$
of an open of $X$. Then $H^p(U, \mathcal{F}) = 0$ for all $p > 0$
and any open $U \subset X$.
\end{lemma}

\begin{proof}
Let $\mathcal{F}$ be a sheaf satisfying the assumption of the lemma.
We will indicate this by saying ``$\mathcal{F}$ has vanishing higher
{\v C}ech cohomology for any open covering''.
Choose an embedding $\mathcal{F} \to \mathcal{I}$ into an
injective $\mathcal{O}_X$-module.
By Lemma \ref{lemma-injective-trivial-cech} $\mathcal{I}$ has vanishing higher
{\v C}ech cohomology for any open covering.
Let $\mathcal{Q} = \mathcal{I}/\mathcal{F}$
so that we have a short exact sequence
$$
0 \to \mathcal{F} \to \mathcal{I} \to \mathcal{Q} \to 0.
$$
By Lemma \ref{lemma-ses-cech-h1} and our assumptions
this sequence is actually exact as a sequence of presheaves!
In particular we have a long exact sequence of {\v C}ech cohomology
groups for any open covering $\mathcal{U}$, see
Lemma \ref{lemma-cech-cohomology-delta-functor-presheaves}
for example. This implies that $\mathcal{Q}$ is also an $\mathcal{O}_X$-module
with vanishing higher {\v C}ech cohomology for all open coverings.

\medskip\noindent
Next, we look at the long exact cohomology sequence
$$
\xymatrix{
0 \ar[r] &
H^0(U, \mathcal{F}) \ar[r] &
H^0(U, \mathcal{I}) \ar[r] &
H^0(U, \mathcal{Q}) \ar[lld] \\
&
H^1(U, \mathcal{F}) \ar[r] &
H^1(U, \mathcal{I}) \ar[r] &
H^1(U, \mathcal{Q}) \ar[lld] \\
&
\ldots & \ldots & \ldots \\
}
$$
for any open $U \subset X$. Since $\mathcal{I}$ is injective we
have $H^n(U, \mathcal{I}) = 0$ for $n > 0$ (see
Derived Categories, Lemma \ref{derived-lemma-higher-derived-functors}).
By the above we see that $H^0(U, \mathcal{I}) \to H^0(U, \mathcal{Q})$
is surjective and hence $H^1(U, \mathcal{F}) = 0$.
Since $\mathcal{F}$ was an arbitrary $\mathcal{O}_X$-module with
vanishing higher {\v C}ech cohomology we conclude that also
$H^1(U, \mathcal{Q}) = 0$ since $\mathcal{Q}$ is another of these
sheaves (see above). By the long exact sequence this in turn implies
that $H^2(U, \mathcal{F}) = 0$. And so on and so forth.
\end{proof}

\end{document}
